\section{Inverse y U-shaped Scaling: cómo se interfieren las subtareas ocultas}

La sección anterior trató la «aparición tardía»; esta sección trata el patrón «primero empeora y luego mejora». Este tipo de curva es la que con más facilidad se malinterpreta como una falla del entrenamiento; el ejemplo de Jason muestra que la tarea completa puede depender de la competencia entre varias capacidades. Un modelo mediano aprendió a corregir citas célebres, pero todavía no aprende a obedecer ``repeat after me'', por lo que obtiene peores resultados que un modelo más pequeño; un modelo más grande, que ya tiene ambas capacidades, vuelve a responder correctamente.

\subsection{Detrás de una curva puede haber varias estrategias}

Esta sección comienza por el fenómeno. Inverse scaling se refiere a que la puntuación de una tarea disminuye al aumentar la escala; U-shaped scaling, a que la puntuación primero baja y después vuelve a subir. La forma de la curva por sí sola no nos dice si se trata de una degradación de la optimización, de un conflicto entre los datos o de una competencia entre subtareas; por eso hay que construir un prompt interpretable que separe una por una las posibles estrategias.

\lecturefigure{jason-slide-016.jpg}{Intuición cuatro: elegir las tareas con cuidado produce inverse o U-shaped scaling}{Jason Wei official deck, p.~16}

\begin{knowledgebox}{Términos clave: inverse y U-shaped}
\begin{tabularx}{\textwidth}{L{0.23\textwidth} X}
\toprule
Término & Significado en el curso \\
\midrule
Inverse scaling & Al aumentar la escala del modelo o el cómputo, la metric de una tarea bien definida disminuye en lugar de aumentar. \\
U-shaped scaling & La misma metric primero baja y luego sube conforme crece la escala, lo que indica que distintas capacidades o estrategias dominan en distintos intervalos. \\
Task selection & Elegir prompts que permitan distinguir entre estrategias en competencia, en lugar de escoger solo el benchmark con el mayor desempeño promedio. \\
\bottomrule
\end{tabularx}
\end{knowledgebox}

\lecturefigure{jason-slide-017.jpg}{El modelo pequeño repite la palabra errónea, el mediano corrige la cita por su cuenta y el grande vuelve a seguir la instrucción}{Jason Wei official deck, p.~17}

Al leer la figura, primero hay que distinguir ``glib'' de ``gold''. El objetivo no es completar la cita célebre, sino obedecer la petición del usuario y repetir la oración escrita mal a propósito. El modelo pequeño la repite «por casualidad», porque no conoce la cita; el modelo mediano conoce la cita pero ignora la instruction; el modelo grande reconoce la cita y además es capaz de anteponer la instruction a su tendencia predeterminada de completar.

\subsection{Descomponer la tarea completa en repeat, repair e instruction following}

La página anterior presentó tres comportamientos; esta sección los formula explícitamente como subtareas. El modelo producirá ``glib'' de manera estable solo cuando sea capaz de repetir el texto, detectar el error en la cita y determinar la prioridad de la instrucción actual. Esto indica que la accuracy final es una función de una combinación de capacidades, y no una medición monótona de un único «grado de inteligencia».

\lecturefigure{jason-slide-018.jpg}{Las scaling curves de tres subtareas ocultas explican la forma de U del conjunto}{Jason Wei official deck, p.~18}

El comportamiento final puede simplificarse como
\[
\hat{y}=f\bigl(R(C),Q(C),I(C)\bigr),
\]
donde $R(C)$ representa la capacidad de repetir, $Q(C)$ la tendencia a corregir citas conocidas, $I(C)$ la capacidad de seguir instrucciones explícitas y $f$ la combinación de decisión de las tres. Un modelo pequeño puede tener $R$ alto, $Q$ bajo e $I$ bajo; un modelo mediano, $R$ alto, $Q$ alto e $I$ bajo; un modelo grande tiene las tres altas, y es $I$ la que determina que no corrija el texto del usuario.

\begin{warningbox}{No interpretes automáticamente la forma de U como «el modelo primero aprende mal y luego aprende bien»}
Una caída en la puntuación puede deberse a que una capacidad nueva se impone sobre una estrategia anterior, y no a una degradación general de las representaciones subyacentes. Un diagnóstico correcto requiere evaluar cada subcapacidad por separado y revisar qué objetivo del prompt debería tener prioridad. Una sola accuracy reduce a la misma etiqueta de correcto/incorrecto tanto «conoce la cita pero no obedece la instrucción» como «no conoce la cita en absoluto».
\end{warningbox}

\teachervoice{Entre 00:21:21 y 00:23:56, Jason dibuja una por una las curvas de las tres subtareas y luego regresa a la tarea completa para explicar la forma de U. Este es el procedimiento de análisis más reutilizable de toda la conferencia: ante una curva extraña, primero hay que preguntarse qué estrategias mezcla la metric.}

\subsection{Práctica de investigación: nunca grafiques un solo punto}

Una vez descompuestas las subtareas, el último paso es colocar también las intervenciones experimentales sobre el eje de escala. Que un fine-tuning, un filtrado de datos o un método de prompt supere a la baseline en un punto de presupuesto no demuestra que siga siendo eficaz al aumentar los datos. Se necesitan al menos varios puntos para determinar si la curva se satura, mantiene su pendiente o se separa cada vez más rápido; además, hay que repetir los experimentos con la misma definición de datos y el mismo protocolo de evaluación, para no confundir con una ventaja del método la fluctuación de un solo punto causada por la inicialización aleatoria o la selección de muestras.

\lecturefigure{jason-slide-019.jpg}{Recomendación general de investigación: graficar scaling curves en lugar de reportar solo la mejora en un punto}{Jason Wei official deck, p.~19}

\begin{lstlisting}[caption={Comparación de las scaling curves de una intervención de investigación y de la baseline}]
budgets = [0.25, 0.5, 1.0, 2.0]
for budget in budgets:
    baseline = train(config="baseline", budget=budget)
    method = train(config="intervention", budget=budget)
    log(budget, evaluate(baseline), evaluate(method))
fit_curve_and_report_uncertainty()
\end{lstlisting}

\begin{importantbox}{Una scaling curve responde al menos tres preguntas}
\begin{enumerate}
  \item \textbf{Data efficiency}: ¿el método solo adelanta el mismo desempeño a un presupuesto menor?
  \item \textbf{Asymptote}: ¿la ventaja se satura pronto, de modo que invertir más ya no la amplía?
  \item \textbf{Slope}: ¿el método cambia la pendiente de crecimiento, de modo que el beneficio es mayor a mayor escala?
\end{enumerate}
Un resultado de un solo punto no permite distinguir entre estos tres casos.
\end{importantbox}

\teachervoice{La conclusión de 00:23:59 a 00:25:51 es muy sencilla: ``just plot scaling curves''. No se pide que cada estudiante entrene un frontier model, sino que haga experimentos escalonados dentro del rango de presupuesto que pueda costear, para no tomar un punto fortuito por una tendencia.}

\subsection{Resumen de la sección}

Inverse y U-shaped scaling suelen indicar que detrás de la metric hay estrategias en competencia. Mediante prompts interpretables, evaluaciones de subcapacidades y curvas con varios presupuestos, los investigadores pueden reformular un «fenómeno extraño» como una hipótesis verificable. Aquí se cierra el ciclo completo de la primera conferencia de Jason: observar los datos, identificar las tareas latentes, descomponer la loss, graficar scaling curves y señalar con claridad la metric y los límites de la evidencia.
